\subsection{Primary comparisons and incremental benefit}
Table~\ref{tab:main} reports all scheduled events; Fig.~\ref{fig:performance}
adds paired block intervals. RQ1 has a mixed answer. Synthetic diffusion S
has lower recall than SB by -0.099 [-0.174, -0.034]; its
IoU difference is -0.010 [-0.017, -0.003].
Intel diffusion S/SB recall is 0.213/0.481, but their background
exceedances are 0.218/0.372.
PEMS bootstrap S appears stronger than SB at nominal alpha
(0.528 versus 0.005 recall), with background rates
0.301 versus 0.000.
Under the descriptive test-control cap .10, corresponding recalls become
0.069 and 0.264, at attainable rates
0.026 and 0.064.
Thus a nominal-threshold ranking is not an isolated test of feature information.
The retrospective analysis does not validate a deployable threshold.

For RQ3, synthetic bootstrap T/TB recall is 0.401/0.532,
with paired difference -0.131 [-0.313, +0.054]; unconditional IoU is
0.093/0.239.
Intel bootstrap T instead exceeds TB by +0.389 [+0.264, +0.523]
recall, but their achieved background rates differ. PE and SE are not interchangeable:
Intel bootstrap PE/SE recall is 0.310/0.380,
with IoU 0.022/0.073.
Switching Intel T to diffusion yields recall 0.060 and background
exceedance 0.288. PEMS temporal contrasts have wide
block intervals. These results do not establish a universal entropy ordering.

For RQ4, synthetic bootstrap ST improves on S by
+0.123 [+0.028, +0.230] recall, but ST minus T is
+0.008 [-0.085, +0.127]. Their primary localization rules differ;
fixed-rule alternatives are retained to distinguish feature and ranking effects.
The direct incremental comparison BST minus B is
+0.003 [+0.000, +0.008] recall for synthetic bootstrap.
The other five dataset/reference combinations have zero mean recall change at
this operating point. This is not equivalence: different feature information
can be suppressed by a maximum and its recalibrated threshold.
Synthetic bootstrap BST minus B IoU is +0.008 [+0.000, +0.017].
Both use the same combined localization rule. A saved maximum-dominance audit
shows that conventional components usually determine the final scan maximum.
The retained causal GDN adaptation has synthetic/Intel/PEMS recall
0.514/0.139/0.157, with background rates
0.135/0.436/0.038.
PEMS bootstrap raw residuals reach 0.264 recall at
0.077 background exceedance. These additional
detectors retain their disclosed support and information-access differences.
For bootstrap T, conditional versus unconditional IoU is
0.231/0.093 on synthetic,
0.195/0.076 on Intel and
0.038/0.015 on PEMS.
Mean best-overlap candidate-group IoU is
0.284/0.413/0.325,
an oracle library diagnostic rather than an upper bound on every sensor ranking.
\begin{table*}[t]\centering\small
\caption{Operational extension results at nominal $\alpha=.10$. P/R: event precision/recall; AP: event-PR envelope area; IoU includes missed events as zero. These are not matched achieved background rates. PE/SE, all comparators and common-support results are retained in the artifact.}\label{tab:main}
\begin{tabular}{llrrrrrrrrrrrr}\toprule
&&\multicolumn{4}{c}{Synthetic (three sizes)}&\multicolumn{4}{c}{Intel Lab}&\multicolumn{4}{c}{PEMS-BAY}\\
Reference&Family&P&R&AP&IoU&P&R&AP&IoU&P&R&AP&IoU\\\midrule
Block & S & 0.256 & 0.285 & 0.220 & 0.024 & 0.251 & 0.208 & 0.083 & 0.017 & 0.211 & 0.528 & 0.179 & 0.035 \\
Block & T & 0.329 & 0.401 & 0.290 & 0.093 & 0.329 & 0.389 & 0.172 & 0.076 & 0.291 & 0.389 & 0.229 & 0.015 \\
Block & ST & 0.308 & 0.409 & 0.272 & 0.075 & 0.329 & 0.389 & 0.172 & 0.070 & 0.291 & 0.389 & 0.229 & 0.018 \\
Block & B & 0.484 & 0.650 & 0.452 & 0.193 & 0.000 & 0.000 & 0.139 & 0.000 & 0.224 & 0.148 & 0.207 & 0.030 \\
Block & BST & 0.486 & 0.653 & 0.456 & 0.200 & 0.000 & 0.000 & 0.139 & 0.000 & 0.224 & 0.148 & 0.207 & 0.030 \\
\midrule
Diffusion & S & 0.197 & 0.170 & 0.196 & 0.010 & 0.158 & 0.213 & 0.086 & 0.014 & 0.373 & 0.287 & 0.281 & 0.006 \\
Diffusion & T & 0.227 & 0.318 & 0.209 & 0.044 & 0.048 & 0.060 & 0.051 & 0.002 & 0.276 & 0.403 & 0.225 & 0.011 \\
Diffusion & ST & 0.197 & 0.170 & 0.196 & 0.009 & 0.048 & 0.060 & 0.051 & 0.002 & 0.376 & 0.287 & 0.300 & 0.001 \\
Diffusion & B & 0.327 & 0.392 & 0.284 & 0.102 & 0.444 & 0.056 & 0.120 & 0.015 & 0.267 & 0.111 & 0.211 & 0.022 \\
Diffusion & BST & 0.327 & 0.392 & 0.284 & 0.103 & 0.444 & 0.056 & 0.120 & 0.011 & 0.267 & 0.111 & 0.211 & 0.022 \\
\midrule
\bottomrule\end{tabular}\end{table*}

\begin{figure*}[!tb]\centering\includegraphics[width=\textwidth]{performance.pdf}
\caption{Operational performance and unadjusted paired 95\% whole-block intervals. IoU assigns zero to missed events. Equal nominal alpha does not imply equal achieved background rates; four real recording blocks give limited precision.}\label{fig:performance}\end{figure*}

\subsection{Information captured and mechanism dependence}
RQ2 is affirmative as an information distinction, not a superiority claim.
In 64 aligned simulations, joint row permutation changes the evaluated
correlation matrix by at most $1.3\times10^{-15}$ (rounded upward).
Mean PE changes from 0.964 to 0.989,
SE from 0.586 to 1.115,
and ACF1 from 0.845 to -0.016; spatial entropy is unchanged.
Conversely, copying lowers mean spatial entropy from 0.877
to 0.134, while mean PE changes only from
0.964 to 0.965.
Offsets, gains, sign changes, quantization, replay, periodic interference,
irregularity, legitimate transitions and dropout remain separate diagnostic cases.
\emph{These deliberately constructed experiments are not the primary benchmark.}

In the heterogeneous synthetic benchmark, bootstrap T/TB recall on flatlining
is 0.741/0.241;
for independent noise it is 0.463/1.000.
These are mechanism-specific operational results, not calibrated guarantees at
identical false-alert rates. A rolling window crossing a flatline onset can
remain measurable; a wholly constant window correctly abstains.
The artifact reports paired mechanism, severity, duration, measured-direction
and affected-support strata. Direction is the measured lagged change, not a
fault-name label or deviation from reference expectation. The first scheduled
in-event decision is already 11 samples late. Short events can therefore be
missed before adequate window evidence becomes available.

\begin{figure*}[!tb]\centering\includegraphics[width=\textwidth]{traces.pdf}
\caption{First prespecified synthetic copy episode: common forecast draws supply all bands. The first affected sensor and the best-overlap candidate group are chosen for illustration only, using injection labels; they are not a detection/localization rule. Shading marks the 12-row injection.}\label{fig:traces}\end{figure*}

\subsection{Reference fidelity, support and calibration}
For RQ5, the shared-support audit in Table~\ref{tab:fidelity} distinguishes
reference defects from feature limits. Synthetic64 diffusion spatial-entropy
coverage is 0.453
versus 0.740 for bootstrap;
Intel values are 0.083 and
1.000, on only
6.000 common spatial units per audit issuance on average.
Intel diffusion has better raw energy score
(6.336 versus 7.234),
yet its sample-entropy coverage is 0.000
versus 0.953 for bootstrap.
Good raw marginal behavior does not ensure a faithful temporal functional.
No consistent generative-reference improvement is established, and this compact
long-lead model does not represent all generative approaches.
\begin{table*}[t]\centering\small
\caption{Paired fidelity on unchanged backgrounds: seed 17, first two blocks, decisions 167/263. References share evaluated raw coordinates or eligible feature units; donor masks can still change per-draw support. Raw, H, PE, SE: 90\% coverage; TV: ordinal-pattern total variation; ES: normalized energy score (lower is better). Temporal entries use W=96. Means omit undefined units; counts and widths are saved.}\label{tab:fidelity}
\begin{tabular}{llrrrrrrrr}\toprule Dataset&Reference&Raw&H&PE&SE&ACF1 MAE&Ordinal TV&R RMSE&ES\\\midrule
Synthetic & Block & 0.875 & 0.812 & 0.855 & 0.858 & 0.034 & 0.081 & 0.236 & 0.578 \\
Synthetic & Diffusion & 0.803 & 0.513 & 0.429 & 0.222 & 0.047 & 0.150 & 0.313 & 1.747 \\
Intel Lab & Block & 0.616 & 1.000 & 0.887 & 0.953 & 0.013 & 0.378 & 0.467 & 7.234 \\
Intel Lab & Diffusion & 0.714 & 0.083 & 0.214 & 0.000 & 0.114 & 0.406 & 0.525 & 6.336 \\
PEMS-BAY & Block & 0.847 & 0.819 & 0.873 & 0.862 & 0.121 & 0.094 & 0.288 & 1.501 \\
PEMS-BAY & Diffusion & 0.739 & 0.494 & 0.699 & 0.430 & 0.163 & 0.118 & 0.360 & 2.055 \\
\bottomrule\end{tabular}\end{table*}

RQ6 exposes distinct operational and statistical limitations (Fig.~\ref{fig:availability}).
Intel S group/window eligibility is 0.017 under bootstrap
and 0.049 under diffusion; T eligibility is
0.541 and 0.633.
Per-sensor support avoids simultaneous complete rows across a group, but
availability alone does not guarantee calibrated detection. Recalibrating
on common feature support can also change achieved rates: Intel bootstrap
common-support B background exceedance is 0.551,
compared with 0.000 operationally. Varying support and nonexchangeable
calibration/test distributions remain unresolved by nominal rank calibration.

Development selects $q=3$ throughout, delay 1 for synthetic and 2 for both real
datasets; $q=4$ loses support at these short windows. Sample tolerance and group
fraction are saved per configuration. At W96, scale 2 can retain sufficient
templates, whereas scale 4 cannot meet the declared 30-template minimum;
W24 is unsupported for the primary temporal entropies. We report these as
limited multiscale support diagnostics; a separately frozen scale-2 detector
sensitivity is reported below.

\begin{figure*}[!tb]\centering\includegraphics[width=\textwidth]{availability.pdf}
\caption{Top: operational eligible-group fraction. Middle: untouched-background exceedance at nominal alpha .10. Bottom: development-only support after independent added missingness, W96; spatial is group support, PE/SE per-sensor support. These different denominators are intentional and explicit.}\label{fig:availability}\end{figure*}

The seed-17, separately calibrated W48/W96 and 51-versus-52-reference-draw
support sensitivity completed 5 of five configurations.
Across completed configurations, the largest absolute recall change caused
by the support-rounding correction is 0.000. Window-specific thresholds
and results are retained separately; the 120-step forecast lead is held fixed.

The separately frozen, exploratory seed-17 scale-2 comparison keeps W96's
physical history, d=24 and both references. All temporal measurements use the
same two-row averages; spatial measurements remain on the original grid.
Bootstrap T recall changes (scale 2 minus 1, paired block intervals) are
-0.093 [-0.231, +0.069], -0.069 [-0.458, +0.333] and
-0.333 [-0.722, -0.028] for synthetic, Intel and PEMS. Corresponding TB
changes are -0.074 [-0.245, +0.042], +0.264 [+0.000, +0.750] and
-0.153 [-0.278, -0.028]. Full PE/SE, localization, support and achieved
background rates accompany these separately calibrated comparisons in the
artifact. Scale-1 calibration ranks reproduce the independent W96 audit;
S/SB scores are unchanged across scales. This reused-data sensitivity does
not select a new primary scale.

A post-hoc CPU estimator check uses 128 independent unit-variance
Gaussian AR(1) realizations. For independent continuous observations, PE's
between-realization SD is 0.023 at W48 and 0.009 at W96.
With fixed sample-entropy tolerance .2, finite supported fractions are
0.086 and 1.000; conditional means therefore require caution.
At $\rho=.8$, quantization to .5-unit steps yields tied-template fraction
0.653. Stable-tie PE eligibility is
0.008, versus 0.758 after discarding ties.
This demonstrates sensitivity of the quality/support decision, not a reason
to reinterpret arbitrary tied ordering as confident dynamics.
